\documentclass[aspectratio=169]{beamer}
\usepackage{fontspec}
\setsansfont{Arial}
\usepackage[spanish]{babel}
\usepackage{graphicx}
\usetheme{Madrid}
\definecolor{uasBlue}{HTML}{123B6D}
\definecolor{uasGold}{HTML}{A77A16}
\setbeamercolor{structure}{fg=uasBlue}
\setbeamercolor{alerted text}{fg=uasGold}
\setbeamertemplate{navigation symbols}{}
\title[Derecho mercantil]{Título de la exposición por completar}
\subtitle{Derecho mercantil}
\author{Estudiante por confirmar}
\institute[UAS-FCA]{Universidad Autónoma de Sinaloa\\Facultad de Contaduría y Administración\\Licenciatura en Contaduría Pública}
\date{Fecha por confirmar}
\titlegraphic{\includegraphics[height=1.1cm]{UAS/licenciatura-en-contaduria-uas/derecho-mercantil/assets-derecho-mercantil/escudo-uas.png}\hspace{1cm}\includegraphics[height=.9cm]{UAS/licenciatura-en-contaduria-uas/derecho-mercantil/assets-derecho-mercantil/fcauas-logo.jpg}}
\begin{document}
\begin{frame}
\titlepage
\end{frame}
\begin{frame}{Problema y propósito}
\begin{itemize}
\item Delimitar el tema jurídico y el propósito de la exposición.
\item Identificar conceptos, normas y fuentes verificadas.
\item Explicar su relación con la Contaduría Pública.
\end{itemize}
\end{frame}
\begin{frame}{Análisis y aplicación}
Sustituir este espacio por un ejemplo o un caso pertinente. Distinguir hechos, norma aplicable, interpretación y conclusión; citar las fuentes realmente consultadas.
\end{frame}
\begin{frame}{Conclusiones y fuentes}
Sintetizar el argumento y sus límites. Incorporar referencias completas y verificar que correspondan a las citas de las diapositivas. Esta plantilla no presupone una exposición obligatoria adicional a las consignas del curso.
\end{frame}
\end{document}